\section{Conclusion and Discussion}
\label{sec:discuss}
We introduce a generalized double descent hypothesis: models and training procedures exhibit atypical behavior
when their Effective Model Complexity
is comparable to the number of train samples.
% This hypothesis significantly generalizes prior work in double descent,
% using a novel notion of model complexity.
We provide extensive evidence for our hypothesis
in modern deep learning settings,
and show that it is robust to choices of dataset,
architecture, and training procedures.
In particular, we demonstrate ``model-wise double descent''
for modern deep networks and
characterize the regime where bigger models can perform worse.
We also demonstrate ``epoch-wise double descent,''
which, to the best of our knowledge, has not been previously proposed.
Finally, we show that the double descent phenomenon can
lead to a regime where training on more data leads to worse test performance.
%All of the above phenomenon follow as natural consequences of our generalized double descent hypothesis.
Preliminary results suggest that double descent also holds
as we vary the amount of regularization for a fixed model
(see Figure~\ref{fig:regular_ocean}).

We also believe our characterization of the critical regime provides a useful
way of thinking for practitioners---if a model and training procedure are
just barely able to fit the train set, then small changes
to the model or training procedure may yield unexpected behavior
(e.g. making the model slightly larger or smaller, changing regularization, etc. may hurt test performance).

\paragraph{Early stopping.} 
We note that many of the phenomena that we highlight
often do not occur with optimal early-stopping.
% That is, when the training procedure keeps
% a validation set, and returns the model at the train-epoch with best validation error.
However, this is consistent with our
generalized double descent hypothesis: 
if early stopping prevents models from reaching $0$ train error
then we would not expect to see double-descent, since 
%we are not allowing the EMC to reach the number of train samples.
the EMC does not reach the number of train samples.
Further, we show at least one setting where model-wise double descent can still occur
even with optimal early stopping
(ResNets on CIFAR-100 with no label noise, see Figure~\ref{fig:app-cifar100-clean}).
We have not observed settings where more data hurts
when optimal early-stopping is used.
However, we are not aware of reasons which preclude this from occurring.
We leave fully understanding the optimal early stopping behavior of double descent
as an important open question for future work.



% In our framework, optimal early-stopping
% can be interpreted as searching for the
% optimal Effective Model Complexity among the 
% then increasing model size does not correspond to
% strictly increasing Effective Model Complexity (EMC)---

\paragraph{Label Noise.}
In our experiments, we observe double descent most strongly in settings with label noise.
However, we believe this effect is not fundamentally about label noise,
but rather about \emph{model mis-specification}.
For example, consider a setting where the label noise is not truly random,
but rather pseudorandom (with respect to the family of classifiers being trained).
In this setting, the performance of the Bayes optimal classifier would not change
(since the pseudorandom noise is deterministic, and invertible),
but we would observe an identical double descent as with truly random label noise.
Thus, we view adding label noise as merely a proxy for making distributions ``harder''---
i.e. increasing the amount of model mis-specification.

\paragraph{Other Notions of Model Complexity.}
Our notion of \emph{Effective Model Complexity} is related to classical complexity notions
such as Rademacher complexity, but differs in several crucial ways:
(1) EMC depends on the \emph{true labels} of the data distribution,
and (2) EMC depends on the training procedure, not just the model architecture.

Other notions of model complexity which do not incorporate features (1) and (2)
would not suffice to characterize the location of the double-descent peak.
Rademacher complexity, for example, is determined by the
ability of a model architecture to fit a randomly-labeled train set.
But Rademacher complexity and VC dimension are both insufficient to 
determine the model-wise double descent peak location, since they do
not depend on the distribution of labels--- and our experiments show that adding label noise
shifts the location of the peak.

Moreover, both Rademacher complexity and VC dimension
depend only on the model family and data distribution, and not
on the training procedure used to find models.
Thus, they are not capable of capturing train-time double-descent effects,
such as ``epoch-wise'' double descent, and the effect of data-augmentation on
the peak location.